\subsection{Real-energy decision outcome}
All nine family-level selected certificates were non-vacuous. The predeclared deployment rule selected a four-rule radius-controlled TSK model for Zone~1 and an eight-rule fixed-$K$ TSK model for Zone~2. No candidate was eligible in Zone~3 because the certification target-clipping rate was 8.01\%, above the predeclared 5\% threshold. The two deployed learned models did not improve the independent test outcome: relative to the 24-hour seasonal-naive fallback, test RMSE changed by $-11.86\%$ in Zone~1 and $-1.51\%$ in Zone~2, while the asymmetric operational cost changed by $-22.67\%$ and $-19.05\%$, respectively. Negative values denote deterioration. Across all three zones, including the fallback decision in Zone~3, mean changes were $-4.46\%$ in RMSE and $-13.91\%$ in cost.

Ridge regression had the tightest mean certificate (0.0790), compared with 0.1462 for radius-controlled TSK and 0.1625 for fixed-$K$ TSK. It also had the best mean test RMSE and cost. Radius control reduced the selected rule count relative to fixed-$K$ in Zones~1 and~2 and tightened the corresponding certificates, but did not improve their test RMSE. This case therefore demonstrates a genuine use of a non-vacuous certificate in a frozen decision protocol, but not an operational success claim. In particular, the certificate controls the clipped sequential squared risk; it does not guarantee dominance in raw RMSE or in an asymmetric downstream cost under later distribution change.
